\documentclass{article}
\usepackage{pseudocode}

\begin{document}
\section{TAD TableDeCodage}

\begin{algorithme}
\signatureFonction{obtenirTableDeCodage}
	{a : ArbreHuffman}{TableDeCodage}
	{non estVide(a)}
\end{algorithme}

\begin{algorithme}
	\signatureFonction{obtenirCode}
		{tab : TableDeCodage, car : \caractere}{CodeBinaire}
		{estPresent(tab,car)}
\end{algorithme}

\begin{algorithme}
	\signatureFonction{obtenirCarac}
		{tab : TableDeCodage, code : CodeBinaire}{\caractere}
		{}
\end{algorithme}

\begin{algorithme}
	\signatureProcedure{ajouter}
		{\paramEntreeSortie{ tab : TableDeCodage} \paramEntree{ car : \caractere, code : CodeBinaire}}
		{}
\end{algorithme}

\begin{algorithme}
	\signatureProcedure{retirer}
		{\paramEntreeSortie{ tab : TableDeCodage} \paramEntree{ car : \caractere}}
		{}
\end{algorithme}

\begin{algorithme}
	\signatureFonction{estPresent}
		{tab : TableDeCodage, car : \caractere}{\booleen}
		{}
\end{algorithme}

\end{document}

